\chapter{Durchführung} \label{ch:Durchfuehrung}

\section{Messverfahren}
\label{sec:Messverfahren}
In der ersten Aufgabe wird das Rauschspektrum eines ohmschen Widerstands qualitativ untersucht. Dazu wird das Gehäuse mit den umschaltbaren Widerständen direkt auf den Verstärkereingang gesteckt, ohne zusätzliche Kabel zu verwenden, um parasitäre Effekte zu minimieren. Der Verstärkerausgang wird an Kanal 1 eines Oszilloskops angeschlossen. Zunächst wird im Oszilloskop-Modus der zeitliche Verlauf der Rauschspannung für verschiedene Widerstandswerte beobachtet. Anschließend wird auf den Spektrumanalyzer-Modus umgeschaltet, um die Verteilung der Rauschleistung über den Frequenzen zu analysieren. Dabei wird der Frequenzbereich zunächst auf 3MHz erweitert, um den Abfall bei hohen Frequenzen sichtbar zu machen. Abschließend wird der Bandfilter zugeschaltet, um den Effekt der Bandbreitenbegrenzung und die Unterdrückung von Störfrequenzen im Spektrum zu beobachten.

\img{schaltkreis}{Schematischer Versuchsschaltkreis. \cite{skriptPAP22}}

Die zweite Aufgabe umfasst die quantitative Messung der Rauschspannung als Funktion des ohmschen Widerstands. Der Versuchsaufbau wird so verändert, dass der Verstärkerausgang über ein kurzes Kabel mit dem Eingang des Bandfilters und dieser wiederum mit dem Voltmeter verbunden wird. Am PC werden für sechs verschiedene Widerstandswerte im Bereich von 5k$\Omega$ bis 30k$\Omega$ jeweils 125 Einzelmessungen durchgeführt. Aus diesen Daten werden der Mittelwert der Rauschspannung sowie der statistische Fehler des Mittelwerts berechnet. Zusätzlich wird die Zimmertemperatur gemessen. Um das Eigenrauschen des Verstärkers zu bestimmen, wird der Verstärkereingang mittels eines Kurzschlusssteckers kurzgeschlossen und die Messung analog zu den Widerstandsmessungen durchgeführt.

In der dritten Aufgabe wird der Frequenzgang der Messelektronik (Verstärker und Bandfilter) bestimmt. Ein Dämpfungsglied wird direkt auf den Verstärkereingang gesteckt, und der Ausgang des Bandfilters wird an das Oszilloskop angeschlossen. Ein Funktionsgenerator liefert ein Sinussignal mit konstanter Amplitude und variabler Frequenz. Das Messprogramm am PC steuert den Funktionsgenerator und das Oszilloskop an, um für eine Vielzahl von Frequenzen die Ausgangsamplitude zu erfassen. Die Messdaten werden als CSV-Datei exportiert. Aus diesen Daten wird der Frequenzgang $g(f)$ berechnet und numerisch integriert, um die äquivalente Rauschbandbreite $B$ zu bestimmen.

Als freiwillige Zusatzübung kann die Rauschspannung als Funktion der Temperatur gemessen werden. Hierfür steht ein beheizbarer Platinwiderstand (Pt4000) zur Verfügung, dessen Widerstandswert temperaturabhängig ist und zur präzisen Temperaturbestimmung genutzt werden kann. Der Widerstand wird auf verschiedene Temperaturen (50$^\circ$C bis 250$^\circ$C) aufgeheizt. Für jede Temperaturstufe wird zunächst der Widerstandswert und daraus die Temperatur berechnet, anschließend wird die Rauschspannung bei dieser Temperatur gemessen. Dies ermöglicht die Überprüfung der Temperaturabhängigkeit des thermischen Rauschens gemäß der Nyquist-Formel und eine weitere Bestimmung der Boltzmann-Konstante.